\subsection{Statistical Analysis}
\label{sec:power}

The statistical analysis evaluates whether analytical outputs generated under
different agent configurations are systematically preferred by human raters.
Because evaluation is performed through pairwise document comparisons, the
analysis focuses on directional preference frequencies rather than absolute
scores.

\paragraph{Rationale for pairwise comparison.}

Evaluation of open-ended analytical outputs is difficult to perform using
absolute scoring scales because raters may apply scoring criteria
inconsistently across cases. Pairwise comparison mitigates this problem by
asking raters to judge which of two analytical outputs more closely matches
the structure and reasoning expected for the task. Prior work in AI agent
evaluation has shown that comparative judgments are more reliable than
absolute ratings for complex analytical tasks. In this study, pairwise
comparisons are used to obtain directional preference observations that can
be aggregated statistically while preserving blinded evaluation of the
agent configurations.
\paragraph{Pairwise comparison structure.}

For each case, three pairwise comparisons are constructed:

\begin{itemize}
\item Baseline vs Protocol
\item Protocol vs Independent
\item Baseline vs Independent
\end{itemize}

Each comparison consists of two normalized documents presented to a blinded
rater as \textit{Document A} and \textit{Document B}. Raters evaluate the pair
independently across four rubric dimensions:

\begin{itemize}
\item R1 — Analytical clarity
\item R2 — Structural reasoning quality
\item R3 — Diagnostic usefulness
\item R4 — Overall analytical usefulness
\end{itemize}

For each dimension, the rater records one of three responses:

\begin{itemize}
\item Document A closer
\item Document B closer
\item No clear difference
\end{itemize}

These responses generate a set of directional preference observations.

\paragraph{Data representation.}

Each rating observation is recorded as a structured record containing:

\begin{itemize}
\item case\_id
\item document\_A\_id
\item document\_B\_id
\item generation\_condition\_A
\item generation\_condition\_B
\item rater\_id
\item evaluation\_dimension
\item response $\in \{A, B, \text{Tie}\}$
\end{itemize}

\paragraph{Preference aggregation.}

For each comparison type and evaluation dimension, preference counts are
aggregated across raters:

\[
N(A>B), \quad N(B>A), \quad N(\text{Tie})
\]

Tie responses are excluded from directional preference calculations but are
reported descriptively.

The directional preference rate favoring configuration $A$ over configuration
$B$ is defined as

\[
\text{PreferenceRate}(A>B) =
\frac{N(A>B)}{N(A>B) + N(B>A)}.
\]

\paragraph{Hypothesis testing.}

The principal hypothesis tests whether the protocol-guided configuration
produces outputs that are more frequently preferred than baseline outputs.

Let

\[
N(P>B)
\]

denote the number of directional preferences in which the Protocol output is
preferred to the Baseline output, and

\[
N(B>P)
\]

the number in which the Baseline output is preferred.

Under the null hypothesis that both configurations are equally likely to be
preferred,

\[
H_0: \Pr(P>B) = 0.5.
\]

Statistical significance is evaluated using a binomial (sign) test applied to
the directional preference counts:

\[
p = \text{BinomialTest}\big(N(P>B),\; N(P>B)+N(B>P),\; p=0.5\big).
\]

\paragraph{Effect size.}

In addition to significance testing, effect size is reported using the
difference in directional preference rates:

\[
\text{EffectSize} =
\text{PreferenceRate}(P>B) -
\text{PreferenceRate}(B>P).
\]

Positive values indicate that protocol-guided analyses are preferred more
often than baseline analyses.

\paragraph{Aggregate and case-level analysis.}

Preference statistics are computed at two levels:

\begin{itemize}
\item \textbf{Case-level analysis}: preference counts are computed separately
for each case.
\item \textbf{Aggregate analysis}: counts are pooled across all cases to
estimate overall differences between configurations.
\end{itemize}

\paragraph{Sensitivity analysis.}

To evaluate robustness of the Independent condition, the analysis is repeated
using alternate assessor extracts when available. Differences between results
obtained with primary and alternate extracts provide a measure of sensitivity
to assessor source selection.

\paragraph{Inter-rater consistency.}

Agreement across raters is evaluated using standard inter-rater reliability
measures such as Fleiss' $\kappa$ or pairwise agreement rates. This analysis
assesses the stability of preference judgments across evaluators.

A lightweight lexical contamination check was used to detect cross-condition
style leakage prior to normalization and blinded evaluation.